%! TeX root = ../main.tex
\chapter{PROFIL PERUSAHAAN}
\label{chap:profilperusahaan}

\paragraphtextstyle

\section{IDENTITAS PERUSAHAAN}
\label{sec:identitasperusahaan}

Nama perusahaan alat kesehatan tempat penulis melakukan magang adalah di \ldots. Alamat perusahaan alat kesehatan tempat penulis melakukan magang adalah di \ldots.

\section{SEJARAH PERUSAHAAN}
\label{sec:sejarahperusahaan}

XXX adalah perusahaan XXXXXX yang berdiri pada tahun XXXXX. Perusahaan ini didirikan oleh \ldots dan seterusnya.

\section{BIDANG USAHA}
\label{sec:bidangusaha}

XXX secara khusus bergerak di bidang XXXXX.

\section{SUSUNAN ORGANISASI}
\label{sec:susunanorganisasi}

\begin{figure}[H]
  \centering
  \fbox{\parbox[c][45mm][c]{0.85\textwidth}{\centering Struktur Organisasi Perusahaan}}
  \caption{Struktur Organisasi PT Mega Andalan Kalasan}
  \label{fig:strukturorganisasi}
  {\fontsize{10}{12}\selectfont Sumber: Kepala Unit Bagian HRD (xxxxx, 20xx)\par}
\end{figure}

\section{PORTOFOLIO PERUSAHAAN}
\label{sec:portofolioperusahaan}

Perusahaan yang telah berdiri selama xx tahun ini memiliki beragam produk yang telah dihasilkan. Berikut beberapa proyek yang pernah ditangani oleh Artmosphere Interior Design Inc.

\clearpage
\begin{table}[H]
  \centering
  \caption{Portofolio PT Mega Andalan Kalasan}
  \label{tab:portofolio}
  \begin{tabularx}{\textwidth}{|c|c|X|}
    \hline
    No & Tahun & Nama Proyek \\
    \hline
    1 & 2016 & Tempat Tidur Rumah Sakit \\
    \hline
    2 & 2016 & Meja Operasi \\
    \hline
    3 & 2016 & Alat Transportasi Pasien \\
    \hline
  \end{tabularx}

  \vspace{0.5ex}
  {\fontsize{10}{12}\selectfont Sumber: Kepala Unit Bagian HRD (xxxxx, 20xx)\par}
\end{table}
